\chapter{Express Server Setup}

\section{Why Split \texttt{app.js} and \texttt{server.js}}

\whatwhyhow{
A production-grade Express backend separates the \textbf{application definition} (\texttt{app.js}) from the \textbf{server bootstrap} (\texttt{server.js}).
}{
\texttt{app.js} defines what the application \emph{is} — middleware, routes, error handlers — with zero knowledge of ports, sockets, or the database. \texttt{server.js} defines how the application \emph{runs} — it starts the HTTP listener and connects to MongoDB. This separation makes the app testable (you can import \texttt{app} into a test file with tools like Supertest without ever binding a port) and keeps startup concerns out of business logic.
}{
Export the configured Express instance from \texttt{app.js}. Import it in \texttt{server.js}, connect to the database, then call \texttt{app.listen(PORT)}.
}

\subsection{backend/src/app.js}

\begin{lstlisting}[language=JavaScript]
const express = require("express");
const cors = require("cors");
const cookieParser = require("cookie-parser");

const authRoutes = require("./routes/auth.routes");
const userRoutes = require("./routes/user.routes");
const taskRoutes = require("./routes/task.routes");

const app = express();

// ---- Core middleware ----
app.use(express.json());                 // parses JSON request bodies -> req.body
app.use(express.urlencoded({ extended: true }));
app.use(
  cors({
    origin: process.env.CLIENT_URL || "http://localhost:5173",
    credentials: true,                    // allow cookies to cross origin
  })
);
app.use(cookieParser());                  // parses cookies -> req.cookies

// ---- Health check ----
app.get("/api/health", (req, res) => {
  res.status(200).json({ status: "ok" });
});

// ---- Routes ----
app.use("/api/auth", authRoutes);
app.use("/api/users", userRoutes);
app.use("/api/tasks", taskRoutes);

// ---- 404 handler ----
app.use((req, res) => {
  res.status(404).json({ message: "Route not found" });
});

// ---- Centralized error handler ----
app.use((err, req, res, next) => {
  console.error(err.stack);
  res.status(err.statusCode || 500).json({
    message: err.message || "Internal Server Error",
  });
});

module.exports = app;
\end{lstlisting}

\subsection{backend/src/server.js}

\begin{lstlisting}[language=JavaScript]
require("dotenv").config();
const app = require("./app");
const connectDB = require("./config/db");

const PORT = process.env.PORT || 5000;

const startServer = async () => {
  try {
    await connectDB();
    app.listen(PORT, () => {
      console.log(`Server running on port ${PORT}`);
    });
  } catch (error) {
    console.error("Failed to start server:", error.message);
    process.exit(1);
  }
};

startServer();
\end{lstlisting}

\begin{diagrambox}[Startup Sequence]
\begin{verbatim}
node server.js
     |
     v
load .env (dotenv)
     |
     v
require app.js  ---> builds Express app + middleware + routes (no side effects)
     |
     v
connectDB()  ------> await mongoose.connect(MONGO_URI)
     |
     v
app.listen(PORT) --> HTTP server starts accepting requests
\end{verbatim}
\end{diagrambox}

\section{Core Concepts Recap}

\begin{itemize}
  \item \textbf{app} --- the Express application object; a callable request handler you configure once.
  \item \textbf{Middleware} --- functions of shape \texttt{(req, res, next)} that run in order for every matching request; they can inspect/modify \texttt{req}/\texttt{res} or short-circuit with a response.
  \item \textbf{Routes} --- map an HTTP method + path to a handler function.
  \item \textbf{req (Request)} --- incoming data: \texttt{req.body}, \texttt{req.params}, \texttt{req.query}, \texttt{req.cookies}, \texttt{req.headers}.
  \item \textbf{res (Response)} --- outgoing data: \texttt{res.status()}, \texttt{res.json()}, \texttt{res.cookie()}, \texttt{res.send()}.
\end{itemize}

\begin{notebox}[NOTE]
Middleware order matters. \texttt{express.json()} must run \emph{before} any route that reads \texttt{req.body}, otherwise the body will be \texttt{undefined}.
\end{notebox}

\begin{warnbox}[WARNING]
Never call \texttt{app.listen()} inside \texttt{app.js}. Doing so binds a port as a side effect of importing the module, which breaks unit testing and can cause duplicate listeners if the module is required twice.
\end{warnbox}
